\chapter{Conclusion and Future Work}
\label{ch:conclusion}

\section{Conclusion}
The objective of this thesis was to develop a turbulent topology optimization framework in FEniCS that is sufficiently rigorous to support later heat-transfer applications, while remaining honest about the physical and numerical limits of the underlying turbulence model. Across Chapters~\ref{ch:turbulence_modeling} to \ref{ch:benchmark_results}, the work progressed in the correct order: first the turbulence model was assessed, then the FEniCS implementation was verified, then the topology optimization methodology was formulated, and only after that was the design framework subjected to an independent CFD re-analysis strategy.

The first research question asked how accurately the Spalart-Allmaras model predicts curvature-dominated internal turbulent flows and where its limitations appear. Chapter~\ref{ch:turbulence_modeling} answered that question in a reasonably clear way. The SA model is attractive for topology optimization because it is wall-resolved, relatively robust, and much easier to integrate into an adjoint-based framework than a two-equation model. However, its limitations are not incidental. The isotropic Boussinesq closure makes it vulnerable in flows with strong curvature-induced anisotropy and separation. Within the benchmark evidence assembled in this thesis, SA remains credible for moderate curved internal flows, while the risk of model breakdown becomes materially more important as the Dean number approaches the upper range identified in Chapter~\ref{ch:turbulence_modeling}.

The second research question asked whether a custom standard-Galerkin FEniCS implementation of the SA model can converge robustly and achieve quantitative agreement with established CFD solvers. Chapter~\ref{ch:fenics_implementation} answered that question positively. The combination of IPCS-based flow solution, careful wall-distance treatment, and wall-resolved meshing was sufficient to produce a working FEniCS-SA solver that can be used as the state engine of a topology optimization loop.

The third research question required the largest correction to the story of the thesis. It is not answered convincingly by showing that turbulent-optimized shapes differ from laminar-optimized shapes. The stronger and more relevant question is whether the final SA-optimized designs remain credible when they are exported and re-analyzed independently in Ansys Fluent. Chapter~\ref{ch:benchmark_results} was therefore restructured around that logic: FEniCS-SA optimization first, Ansys-SA re-analysis second, and Ansys SST $k$-$\omega$ sensitivity third. That is the correct basis for the final claim of the thesis.

At the present draft stage, the honest overall conclusion is therefore mixed but coherent. The framework itself is credible: the turbulence model has a clearly defined operating window, the FEniCS implementation has been verified, and the topology optimization methodology is mathematically consistent. What remains conditional is the final quantitative statement of Chapter~\ref{ch:benchmark_results}. That chapter now contains the correct figures, tables, and interpretation logic, but its decisive claims must be filled with the completed Ansys campaign before the thesis can state how close FEniCS-SA is to Ansys-SA and how strongly the final designs depend on the turbulence model.

\section{Limitations and Future Work}
The main limitations of the thesis are no longer hidden in the margins. They are central to how the results should be interpreted. The most important ones are the approximate nature of the frozen adjoint, the difference between cross-code verification and true validation, the ambiguity introduced by geometry extraction from a diffuse design field, and the limited turbulence-model hierarchy available within the thesis timescale.

\subsection{Frozen Adjoint as a Tractable Approximation}
The frozen-turbulence adjoint remains the practical optimization engine of the thesis. That choice should be defended precisely. It is not adopted because it is physically exact, and not because ``efficiency'' is a rhetorical convenience. It is adopted because a fully coupled turbulent adjoint is substantially harder to implement, more fragile numerically, and more difficult to stabilize within the scope of a Master's thesis.

This has a direct consequence for the interpretation of the optimized designs. The frozen adjoint can generate useful design directions, but it does not provide exact sensitivities to future turbulence production. Where separation or anisotropic shear dominates, the resulting design may still need a stronger external check. That is why the Ansys re-analysis is not a decorative addition to the thesis. It is the step that makes the optimization claims worth trusting.

\subsection{Cross-Code Verification Is Not Experimental Validation}
Even if the final Chapter~\ref{ch:benchmark_results} campaign shows close agreement between FEniCS-SA and Ansys-SA, that result should still be described carefully. Agreement between two CFD solvers is stronger evidence than an internal self-check, but it is not the same as validation against measurements. The thesis should therefore avoid claiming more than it has demonstrated.

The most defensible final statement is that the FEniCS framework is corroborated by independent commercial CFD and that the model sensitivity of the final designs has been quantified. A later research step, beyond the present thesis, would be to compare at least one extracted geometry against experimental pressure-drop or velocity-profile data. That would move the work from strong numerical corroboration toward genuine physical validation.

\subsection{Geometry Extraction and Benchmark Choice}
The transition from the diffuse topology optimization field to a sharp Ansys geometry remains a real source of uncertainty. Thresholding, contour cleanup, inlet and outlet non-design treatment, and remeshing can all shift the final losses. If that step is under-documented, the Chapter~\ref{ch:benchmark_results} comparison becomes much harder to interpret.

For the final submission, the benchmark choice should also be managed pragmatically. One well-executed curved benchmark is more valuable than two incomplete ones. The preferred route is to complete the Alexandersen/Bayat pipe-bend case first and then add the U-bend only if the additional Ansys campaign can be executed properly. If the U-bend cannot be completed to the same standard, it should remain in future work rather than weakening the core chapter.

\subsection{Model Hierarchy, Three-Dimensionality, and Heat Transfer}
The turbulence-model hierarchy in the thesis is still limited. SA is the optimization model, SST $k$-$\omega$ is the main higher-fidelity comparator, and standard $k$-$\epsilon$ is useful mainly as a secondary reference and for consistency with parts of the literature. This is already enough to say something meaningful about model sensitivity, but it is not the end of the story.

The longer-term extension remains clear. A stronger framework would couple the present methodology to three-dimensional routing problems and later to conjugate heat transfer. However, it would be a mistake to move to that stage before the current two-dimensional CFD re-analysis loop is complete. The immediate next step is not a more ambitious physics package. It is to finish Chapter~\ref{ch:benchmark_results} properly, report the Ansys-SA versus FEniCS-SA difference on the extracted geometries, quantify the SST sensitivity, and state the limits of the SA-driven optimizer without exaggeration.
